\section{A. A dictionary which retains the primes and the extensions}
\label{sec:dictionary}

We use ordinary left matrix multiplication. A prime on a group in the manuscript means a quotient by a specified central subgroup, or an orbifold quotient; it never means a commutator subgroup in this dictionary. We write $F=\langle x,y\rangle$, $z=(xy)^{-1}$, and $\widehat F$ for its full profinite completion. In the notation of \S46 we choose
\[
 x=\ell _1,\qquad y=\ell _0,\qquad z=\ell _\infty,
 \qquad xyz=1.
\]
The switch of the first two labels is essential: the manuscript's product is $\ell _\infty\ell _1\ell _0=1$. A symbol such as $\rho(a)$ below means conjugation by the indicated lift of $\rho$, not multiplication by that lift. Inner automorphisms are $\operatorname{Int}(a)(b)=aba^{-1}$. All automorphisms of profinite groups are continuous.

\begin{longtable}{p{.18\textwidth}p{.40\textwidth}p{.31\textwidth}}
\toprule
Our notation / his notation&Definition and present interpretation&Manuscript location (cote 140-3 unless indicated)\\\midrule\endhead
$F=\Pi_{0,3}$&Tripod fundamental group; discrete free group of rank two. Its completion is not the modular group.&\S42, fols.332--336, scans85--94.\\
$B=\Pi'_{0,3}=\mathfrak S^+_{0,3}=\mathcal G^+_{0,3}$&Oriented modular orbifold group $C_2*C_3=\operatorname{PSL}_2(\mathbb Z)$; $F$ is its index-six level-two kernel.&\S38, fol.297, scan18; \S42, fols.334--336, scans91--94.\\
$B^\pm=\mathfrak S_{0,3}$&Extended modular group $\operatorname{GL}_2(\mathbb Z)/\{\pm I\}$, equivalently the reflection group of signature $(2,3,\infty)$.&\S42, fols.335--337, scans92--96.\\
$\sigma=\sigma_\infty,\rho$&Generators of $B$ with $\sigma^2=\rho^3=1$. On $F$, $\sigma(x,y)=(y,x)$ and $\rho(x,y,z)=(z,x,y)$.&\S42, fol.334$'$, scan91, (33).\\
$J=\tau_\infty$; $\tau$&$\tau$ is the outer complex-conjugation class. $\tau_\infty$ is its genuine lift commuting with $\sigma_\infty$; $J=\operatorname{diag}(1,-1)$ projectively.&\S42, fols.332$'$--335$'$, scans87--93.\\
$\varepsilon_0,\varepsilon_1,\varepsilon_\infty$&$\sigma\rho,\rho\sigma,\rho^{-1}\sigma\rho^{-1}$ in $B$; their squares are $\ell_0,\ell_1,\ell_\infty$.&\S42, fol.333, scan88 (22), and scan92 (34).\\
$\widetilde M=\widetilde M_{0,3}$&Genuine loop automorphisms of $\widehat F$ whose outer class normalizes geometric $S_3$; equivalently the compatible automorphisms of $\widehat B$.&\S46, fol.390, scan182 (1)--(3).\\
$\widetilde N=\widetilde N_{0,3}$&$\widetilde M/\operatorname{Inn}(\widehat F)$; equals $S_3\times C$, where $C$ is the centralizer of $S_3$ in the label-preserving outer loop group.&\S46, fols.390--391, scans182--184.\\
$C=\widetilde N_{0,3}^{!}$&The above centralizer; identified below with $\widehat{\mathrm{GT}}_0$, with the automorphism condition retained.&\S46, scan183 (9bis)--(11).\\
$\widetilde M^!=\widetilde M_{0,3}^{!}$&Inverse image of $C$ in $\widetilde M$; $1\to\widehat F\to\widetilde M^!\to C\to1$.&\S46, scan183 (9bis).\\
$\widetilde\Gamma$; \S45 $\widetilde N_{0,3}$&The quotient by $\widehat B$, hence $C$, before the renaming in \S46. It must not be confused with \S46's $S_3\times C$.&\S45, fol.384, scan170 (1); scans177--178 (20)--(21).\\
$\widetilde N_{1,1}$&Pullback of $\operatorname{GL}_2(\widehat{\mathbb Z})\to\operatorname{GL}_2(\widehat{\mathbb Z})/\{\pm I\}$ along the claimed congruence homomorphism from $\widetilde M$; its geometric kernel is $\widehat{\operatorname{SL}_2(\mathbb Z)}$.&\S45, fols.387--388, scans176--178 (19)--(20).\\
$\widetilde M^![0]$&Normalizer of $\overline{\langle\ell_0\rangle}$; kernel of its map to $C$ is that procyclic group. Round brackets $(0)$ denote a chosen section, after normalization; square brackets do not.&\S46, fols.404--413, scans210--227.\\
$\widetilde M^!(1/2)$, $(2)$, $(-1)$&Centralizers of the respective lifts $\sigma_\infty,\rho\sigma_\infty\rho^{-1},\rho^{-1}\sigma_\infty\rho$; the latter two sections are conjugates of the first.&\S46, fols.402--405$'$, scans206--213.\\
$\widetilde M^!(-\bar j)$&Centralizer of $\rho$ above the subgroup $\mu\equiv1\pmod3$; a partial section, not a section over all $C$.&\S46, fol.403$'$, scan209 (106)--(108).\\
$g,h,f_G;\alpha_0,\beta_0$&$u(\rho)=g\rho$, $u(\sigma)=h\sigma$, $u(\varepsilon_0)=f_G^{-1}\varepsilon_0^\mu f_G$; $g=\alpha_0\rho(\alpha_0)^{-1}$, $h=\beta_0\sigma(\beta_0)^{-1}$.&\S46, fols.392--393, scans186--188.\\
$\alpha,\beta;\xi,\eta$&$\alpha=f_G\alpha_0$, $\beta=f_G\beta_0$; $\xi=\alpha\rho(\alpha)^{-1}$, $\eta=\beta\sigma(\beta)^{-1}$. Simultaneous left multiplication of $\alpha,\beta,f_G$ by $\ell_0^n$ is the gauge ambiguity.&\S46, fols.393--394, scans189--190 (35)--(46).\\
$\zeta$&Not a further independent unknown in the checked \S46 system. In \S44 the readings $\xi\mid\zeta$ refer to a cocycle letter, whose coordinate convention changes by $\underline\xi=\ell_0^\nu\xi\ell_1^{-\nu}$.&\S44, scans151--153; \S46, fols.397--397$'$, scans196--197.\\
$C^\varphi=\widetilde N_{0,3}^{!\approx}$&Classes in $C$ compatible, up to a specified inner conjugacy in $\widehat B$, with $\varphi:\widehat F\twoheadrightarrow\widehat B$, $(\ell_0,\ell_1,\ell_\infty)\mapsto(\sigma\rho,\sigma,\rho^{-1})$.&\S47, fols.415--420, scans229--240.\\
$\widetilde M^{!\approx},\widetilde N^\approx,\widetilde M^\approx$&Inverse images of $C^\varphi$, $S_3\times C^\varphi$, respectively, in the preceding exact sequences. The double mark imposes an additional compatibility.&\S47, fols.420--421, scans239--242 (24)--(39).\\\bottomrule
\end{longtable}

\begin{claim}{A.1}{Refuted}{Refuted; corrected Legendre identification proved}
The literal identifications $M_{1,1}[2]\simeq U_{0,3}$ and $\Pi_{0,3}\simeq\mathfrak S^+_{0,3}$ in the development prompt are false. Over $S=\operatorname{Spec}\mathbb Z[1/2]$ the correct first statement is
\[
 \mathcal M_{1,1}[2]\mathbin{/\mkern-6mu/}\mu_2\simeq
 U=\mathbb P^1_S\setminus\{0,1,\infty\}=\mathcal M_{0,4},
 \qquad \mathcal M_{1,1}[2]\simeq U\times B\mu_2.
\]
Here $/\mkern-6mu/$ denotes rigidification by the central automorphisms $\{\pm1\}$; we use this symbol only in this display. The correct second statement is $\Pi'_{0,3}\simeq\mathfrak S^+_{0,3}$, with a prime. These distinctions are already explicit in the manuscript: cote 140-2, \S37, fols.279--280, scans547--549, and cote 140-3, \S38, fol.297, scan18 (24)--(25).

\end{claim}
\begin{proof}
An ordered full level-two structure labels the three nonzero two-torsion sections. The quotient $E/\{\pm1\}$ is a projective line bundle with four ordered branch sections. Three of them uniquely normalize its coordinate to $0,1,\infty$; the fourth is a section $\lambda$ of $U$. Conversely the double cover
\[
 E_\lambda:\quad v^2=t(t-1)(t-\lambda)
\]
with its point at infinity and its labeled points $(0,0),(1,0),(\lambda,0)$ is an elliptic curve over $U$: its discriminant is $16\lambda^2(1-\lambda)^2$, a unit. Every cover with these four ordered branch sections is an etale-local quadratic twist of this one. Its automorphisms over the ordered branch line are exactly the two deck transformations. Descent therefore identifies the groupoid over $\lambda$ with the groupoid of $\mu_2$-torsors, compatibly with base change. The exhibited family neutralizes the gerbe. Nontrivial inertia $\mu_2$ proves that the unrigidified stack cannot be the scheme $U$. In characteristic two the four-point etale branching construction is unavailable, so extending this identification over $\mathbb Z$ would change the statement.

The tripod retracts onto a bouquet of two circles, hence its topological fundamental group is free of rank two and torsion-free (a nonempty cyclically reduced word has nonempty positive powers). The orbifold obtained by filling two of its punctures with stabilizers of orders two and three has presentation
$\langle\ell'_0,\ell'_1,\ell'_\infty\mid
\ell'_\infty\ell'_1\ell'_0=1,(\ell'_1)^2=(\ell'_\infty)^3=1\rangle$.
Eliminating $\ell'_0$ gives $C_2*C_3$, which contains nontrivial torsion. Thus it cannot be $F$. The exact quotient is $F/\langle\!\langle x^2,z^3\rangle\!\rangle$, after the change of labels above.
\end{proof}

\begin{claim}{A.2}{Proved}{the modular and reflection avatars}
There are isomorphisms of discrete groups
\[
 B=C_2*C_3\simeq\operatorname{PSL}_2(\mathbb Z),\qquad
 B^\pm\simeq\operatorname{GL}_2(\mathbb Z)/\{\pm I\},
 \qquad
 1\longrightarrow F\longrightarrow B\longrightarrow S_3\longrightarrow1.
\]
The last map is reduction modulo two, and the indicated $F$ is free on $\ell_0,\ell_1$. On complex orbifolds, $B$ is the orbifold fundamental group of $[U/S_3]$, or of the rigidified moduli stack of elliptic curves; $B^\pm$ is its extension by complex conjugation. These identifications account for the reflection, cartographic, modular, and universal-surface avatars of \S42, fols.336$'$--337, scans95--96. They induce isomorphisms of full profinite completions, and the index-six exact sequence remains exact after completion.

\end{claim}
\begin{proof}
The matrices
\[
 S=\begin{pmatrix}0&-1\\1&0\end{pmatrix},\quad
 R=\begin{pmatrix}0&1\\-1&1\end{pmatrix},\quad
 J=\begin{pmatrix}1&0\\0&-1\end{pmatrix}
\]
satisfy $S^2=R^3=-I$. The Euclidean algorithm shows that $S$ and $SR=\left(\begin{smallmatrix}1&-1\\0&1\end{smallmatrix}\right)$ generate $\operatorname{SL}_2(\mathbb Z)$. The usual modular triangle in the upper half-plane has one ideal vertex and angles $\pi/2,\pi/3$; reflecting in its sides tiles the half-plane with disjoint interiors. Crossing a side records a reflection, and canceling a backtrack or a complete turn at either finite vertex gives precisely
\[
 r_0^2=r_1^2=r_\infty^2=(r_0r_\infty)^2=(r_1r_0)^3=1.
\]
A reduced path through triangle interiors cannot close without these elementary cancellations: remove an outermost triangle in a finite disk bounded by the path, and induct on the number of triangles. This proves the presentation, rather than only a surjective matrix representation. Its even subgroup is $C_2*C_3$. The reflection matrices can be taken as $r_0=SJ$, $r_1=R^{-1}r_0$, $r_\infty=J$, projectively.

The quotient $U\to U/S_3$ is the map
\[
 q(t)=\frac{27}{4}\frac{t^2(t-1)^2}{(t^2-t+1)^3}.
\]
It has ramification orders two and three over $1,\infty$, and a missing cusp over $0$. Local lifting of paths across a branch point replaces its loop by the corresponding order-two or order-three orbifold generator. This establishes the covering equivalence and the fundamental-group exact sequence. Its kernel is the fundamental group of $U$, hence is $F_2$. Directly, its peripheral generators are
\[
 L_0=(SR)^2=\begin{pmatrix}1&-2\\0&1\end{pmatrix},\quad
 L_1=(RS)^2=\begin{pmatrix}1&0\\2&1\end{pmatrix},\quad
 L_\infty=\begin{pmatrix}3&-2\\2&-1\end{pmatrix}.
\]
Their product $L_\infty L_1L_0=-I$ is the identity in $B$, not in $\operatorname{SL}_2(\mathbb Z)$. Marking a fourth point on the three-punctured sphere identifies its moduli parameter with $U$; allowing permutations of the other three points gives the same quotient orbifold. An oriented marked torus has mapping class group $\operatorname{SL}_2(\mathbb Z)$, as seen by straightening a meridian and longitude and then the complementary disk; quotienting the central elliptic involution gives $B$. Allowing orientation reversal gives $B^\pm$. Thus these interpretations preserve the named peripheral classes.

For a finite-index subgroup $F\subset B$, every finite-index normal subgroup of $F$ contains the core in $B$ of a finite-index subgroup, so the topology induced from $B$ is its full profinite topology. Taking inverse limits gives the claimed exact sequence. No congruence-subgroup assertion is used.
\end{proof}

\begin{claim}{A.3}{Known}{which group is $\widehat{\mathrm{GT}}_0$}
For the \S46 definitions,
\[
 \widetilde N^!\simeq\widehat{\mathrm{GT}}_0,\quad
 \widetilde N\simeq S_3\times\widehat{\mathrm{GT}}_0,\quad
 1\to\widehat F\to\widetilde M^!\to\widehat{\mathrm{GT}}_0\to1,
\]
and $1\to\widehat B\to\widetilde M\to\widehat{\mathrm{GT}}_0\to1$.
In particular, neither genuine-automorphism group $\widetilde M$ nor $\widetilde M^!$ is being identified with the pair group. The identification of the outer centralizer is Harbater and Schneps (2000), Main Theorem (a), or their \S1.3, Theorem 2 \citep{HS2000}; the reduction to their hypotheses follows below.

\end{claim}
% anchor:A.3
\manuscriptanchor{\S 45; fols384--388; scans170--178; \S 46; fols390--391; scans182--184}

\begin{proof}
Let $O$ denote outer loop automorphisms, allowing permutation of the three peripheral classes and a common unit exponent. The geometric $S_3\subset O$ maps isomorphically to the permutation group of these classes. If $v$ normalizes this $S_3$, its conjugation on $S_3$ is inner: an automorphism of $S_3$ permutes its three order-two subgroups, and this embeds $\operatorname{Aut}(S_3)$ into their permutation group, already exhausted by conjugation. Since $Z(S_3)=1$, there is a unique $s\in S_3$ such that $s^{-1}v$ centralizes $S_3$. Hence its normalizer is $S_3\times C_O(S_3)$. An element centralizing $S_3$ has a peripheral permutation in $Z(S_3)$, so preserves labels. Under $\pi_1(\mathcal M_{0,4})=F_2$, this is exactly the quasi-special symmetric group $\operatorname{Out}^{\#}_4$ of Harbater--Schneps: the $S_4$ action on the cross-ratio line factors through $S_4/V_4=S_3$, and all their inertia classes are the three peripheral classes. Their theorem therefore identifies $C_O(S_3)$ with $\widehat{\mathrm{GT}}_0$. Statement B.m2 supplies the explicit coordinate conversion and checks both directions of the equations. Taking inverse images in genuine automorphisms gives the first extension; taking the inverse image of the geometric $S_3$ gives $\widehat B$, by the preceding based orbifold construction and the trivial center of $\widehat F$. Quotienting gives the second extension. This also explains the change of notation between \S45 and \S46.
\end{proof}

\begin{claim}{A.4}{Proved}{finite checks and the P060 alternatives}
The displayed matrix relations have been checked exhaustively at levels $2,3,4,6,12$ (and additionally $5,7,8,24$) by \texttt{code/dictionary\_matrices.py}. The distinction $L_\infty L_1L_0=-I$ persists at every level except two. At level two $S$ and $R$ generate $S_3$, while the image of $F$ generated by $L_0,L_1$ is trivial. Thus calling $\operatorname{SL}_2(\mathbb Z/2)$ a quotient of $F$ via these particular matrices would lose the actual map.

For P060 scan18, the first and second returns give terminal factors $\tau'_\infty\tau'_0$ and $\tau'_0\tau'_\infty$ in the expression for $\ell_1$. Both yield the same element, since $(\tau'_0\tau'_\infty)^2=1$ and both factors are involutions. Mathematics does not choose between those two readings. In the same scan, $(\ell'_1)^2=1$ is correct, whereas $(\ell'_\infty)^2=1$ is false: the latter generator has order three. On scan19 the avatar is a genuine automorphism group, not the outer group (which in the discrete tripod case is finite). Scan20's ``rank one'' is incompatible with $\Im\tau\ne0$; the lattice has rank two. These conclusions are editorial evidence; they do not replace the scan readings.

\end{claim}
% anchor:A.4
\manuscriptanchor{\S 38--39; fols297--298; scans18--20; \S 42; fol336; scan94}

\begin{proof}
The polynomial matrix multiplications used in Statement A.2 establish the identities over $\mathbb Z$, so reduction proves their validity at every finite level. The program independently forms the finite generated groups and tests each identity, including the central signs, by direct multiplication. For the two scan18 words, swapping the last two factors is allowed by their proved commutation, giving equality before any completion. A nonreal $\tau$ makes $1,\tau$ linearly independent over $\mathbb R$, so the integer span is a discrete rank-two subgroup of $\mathbb C$; it cannot have rank one.
\end{proof}

\begin{claim}{A.5}{Known}{All ten avatars, with the fixed fourth puncture retained}
The ten items on \S42, fols.336$'$--337, scans95--96, give the following realizations of $B^\pm$; taking orientation-preserving morphisms gives $B$ in each row. A conformal groupoid in this table permits antiholomorphic morphisms. Ordinary holomorphic orbifold fundamental groups give only the oriented subgroup.
\begin{center}\small
\begin{tabular}{rp{.73\textwidth}}
\toprule
Item&Precise realization of $B^\pm$\\\midrule
1&$\operatorname{Aut}_{\rm lac}(F)$, genuine discrete loop automorphisms.\\
2&Peripheral outer automorphisms of the four-punctured sphere fixing the fourth peripheral class and permuting the first three.\\
3&Components of the homeomorphism group of that sphere fixing the fourth puncture.\\
4&Fundamental group of the quotient groupoid $[U/(S_3\times\langle\text{conjugation}\rangle)]$.\\
5&Extended orbifold group of the sphere with one cusp and orbifold orders two and three.\\
6&The three-reflection cartographic group with the presentation in A.2.\\
7&Fundamental group of the universal conformal tripod, whose total groupoid is the quotient in item 4.\\
8&$\mathcal T_{1,1}/\{\pm1\}$, the torus mapping class group after quotienting its center.\\
9&Fundamental group of the rigidified conformal elliptic moduli groupoid $[\mathbb H/(\operatorname{GL}_2(\mathbb Z)/\{\pm I\})]$.\\
10&$\operatorname{GL}_2(\mathbb Z)/\{\pm I\}$.\\\bottomrule
\end{tabular}\end{center}
The topological input for items 2--3 is the punctured Dehn--Nielsen--Baer theorem and the Birman exact sequence: Farb--Margalit, version 5.0, Theorems 8.8 and 4.6, which we opened and checked in the hyperbolic cases used here. We give the reduction so that a full four-point mapping class group is not substituted for its fourth-point stabilizer.
\end{claim}
\begin{proof}
For item 1, a discrete loop automorphism has a peripheral permutation and a common orientation sign. After removing representatives of these in the reflection group, it preserves the three peripheral conjugacy classes with positive sign. Conjugate to make it fix $x$; write its image of $y$ as $g y g^{-1}$. It preserves the conjugacy class of $xy$. Choose a shortest representative of the double coset $\langle x\rangle g\langle y\rangle$. If nonempty, it begins with a $y$-letter and ends with an $x$-letter. Then $xgyg^{-1}$ is cyclically reduced of length $2+2|g|>2$, contradicting conjugacy to $xy$. Thus $g=x^m y^n$ before minimization, and the automorphism is inner. This proves that the outer group is precisely $S_3\times C_2$. Conjugation from $B^\pm$ has that quotient and geometric kernel $F$. Its kernel is trivial, since a matrix commuting projectively with both $L_0,L_1$ is scalar. Hence it is the required genuine automorphism group.

For items 2--3, Dehn--Nielsen--Baer applies to the hyperbolic four-punctured sphere and respects the peripheral permutation, so restricts to the subgroup fixing the fourth puncture. The Birman sequence for adding that point to the hyperbolic tripod has kernel $\pi_1(U)=F$. Point pushing acts on the tripod group by the corresponding inner automorphism. The resulting homomorphism to item 1 is the identity on $F$ and on the outer quotient $S_3\times C_2$; exactness makes it an isomorphism. This verifies every hypothesis and the precise fixed-label restriction of the imported theorems \citep[Theorems 4.6 and 8.8]{FarbMargalit2011dictionary}.

Items 4--7 are the quotient orbifold, its covering presentation, and its chamber presentation proved in A.2. The universal family of a tripod has no remaining complex moduli, so its total groupoid is the tripod with its conformal symmetry group acting; this identifies item 7 with item 4. For items 8--9, the marked elliptic Teichm\"uller space is the contractible upper half-plane, and change of an integral homology basis acts by the matrices and the conjugation rule of scan92 (37). The orbifold universal cover therefore has deck group $\operatorname{GL}_2(\mathbb Z)$; rigidification removes exactly the central elliptic involution. This gives item 10 and its oriented subgroup, as asserted. None of these comparisons asserts a congruence-completion isomorphism.
\end{proof}
